\section{Numerical Methods}
\label{sec:methods}

\subsection{Two-dimensional weakly damped symplectic lattice}
\label{sec:2d}

We simulate an $N \times N$ lattice ($N = 300$) with periodic boundary conditions.
The integrator is St\"ormer--Verlet with an optional Hubble-like drag $\eta\,\dot\phi$, which is exactly symplectic only when $\eta = 0$ and is weakly area-contracting (conformally symplectic) when $\eta > 0$.
For the 2D Gaussianity ensemble (\cref{sec:gauss}) and the acoustic-peak runs (\cref{sec:peaks}) we set $\eta = 0$, so symplectic structure is preserved exactly up to floating-point error; weak damping ($\eta \in [0.01, 0.06]$) is introduced only in the Phase~I cosmic-timeline run (\cref{sec:timeline}) and in the inverse reconstruction (\cref{sec:inverse}), where it is needed to stabilize gradient flow.
Whenever a quantitative claim in this paper depends on the symplecticity of the integrator we make this distinction explicit; all sections that combine the lattice with the inverse reconstruction or the late-time continuation should be read as ``weakly damped symplectic'', not as ``exactly symplectic''.

The initial field $\phi_0(\bm{x})$ is drawn from a scale-invariant spectrum: in Fourier space, the amplitude is $\hat{\phi}(\bm{k}) \propto |\bm{k}|^{-3/2}$, giving $P(k) \propto k^{-3}$.
This mimics the approximately scale-invariant primordial power spectrum.
The field is normalized to unit variance with zero mean.

Evolution follows \cref{eq:sv} with a 2D isotropic Laplacian kernel:
\begin{equation}
  K = \frac{1}{12}\begin{pmatrix} 1 & 2 & 1 \\ 2 & -12 & 2 \\ 1 & 2 & 1 \end{pmatrix}\,,
  \label{eq:kernel}
\end{equation}
applied via convolution with periodic (wrapped) boundary conditions.
Unless otherwise stated, the baseline parameters are $c^2_{\mathrm{base}} = 0.45$, $c_0 = 10$, $\eta = 0$ (no drag), and the evolution runs for 45 steps.
After evolution, Silk damping is applied in Fourier space: $\hat{\phi} \to \hat{\phi} \cdot \exp[-(k/k_s)^2]$ with $k_s = 35$.

The ``temperature'' field is defined as the normalized fluctuation $T(\bm{x}) = [\phi(\bm{x}) - \bar{\phi}] / \sigma_\phi$.

\subsection{Three-dimensional lattice and sphere projection}
\label{sec:3d}

For full-sky CMB map generation, we extend to a $N^3$ cubic lattice ($N = 96$--$128$) with the 3D Laplacian:
\begin{equation}
  \nabla^2_{\mathrm{lat}} \phi = \frac{1}{6}\sum_{i=1}^{3} (\phi_{\bm{x}+\hat{e}_i} + \phi_{\bm{x}-\hat{e}_i}) - \phi_{\bm{x}}\,,
\end{equation}
and the St\"ormer--Verlet update \cref{eq:sv} with parameters $c^2_{\mathrm{base}} = 0.45$, $\eta = 0.01$--$0.015$, and an optional weak nonlinear term $-\alpha\,\phi^2$ ($\alpha = 0.005$) to model gravitational structure formation.
The initial 3D field uses amplitude spectrum $|\hat{\phi}(\bm{k})| \propto k^{-3/4}$, giving $P(k) \propto k^{-3/2}$ in 3D; this corresponds to a spectral index chosen to produce scale-invariant variance per logarithmic $k$-bin in three dimensions, analogous to the $P(k) \propto k^{-3}$ used in the 2D case.

Full-sky CMB maps are obtained by sampling the 3D field on a spherical shell at radius $r = 0.35 N$ centered on the lattice midpoint.
We use two projection methods:
\begin{enumerate}
  \item \textbf{Mollweide grid sampling}: a $400 \times 800$ grid in $(\theta, \varphi)$ is mapped to Cartesian indices $(x, y, z) = (N/2 + r\cos\theta\cos\varphi, N/2 + r\cos\theta\sin\varphi, N/2 + r\sin\theta)$, with nearest-neighbor lookup.
  \item \textbf{HEALPix sampling}: using \texttt{healpy}~\citep{Gorski2005}, pixel centers at $\nside = 64$ are mapped to 3D coordinates and sampled via trilinear interpolation (\texttt{scipy.ndimage.map\_coordinates}).
\end{enumerate}

\subsection{Dissipative baseline: coupled map lattice}
\label{sec:cml}

For ablation, we compare against a dissipative coupled map lattice (CML)~\citep{Kaneko1992}:
\begin{equation}
  \phi_{n+1}(\bm{x}) = f[\phi_n(\bm{x})] + \varepsilon\,\nabla^2_{\mathrm{lat}} f[\phi_n(\bm{x})]\,,
  \label{eq:cml}
\end{equation}
where $f(x) = 1 - \mu(n) x^2$ is the logistic map with the DSC cooling schedule \cref{eq:cooling} ($\mu_c = 1.5437$, $\kopt = 1.248$), and $\varepsilon = 0.55$ is the diffusion coefficient.
This model is first-order (no $\phi_{n-1}$ term) and dissipative ($|\det DF| < 1$), violating Assumption~2.
We evolve for 150 steps with rescaled initial amplitude ($\sigma_0 = 0.1$) for numerical stability.

\subsection{Analysis pipeline}
\label{sec:analysis}

\paragraph{Power spectrum.}
We compute the radially averaged 2D power spectrum $P(k) = \langle |\hat{\phi}(\bm{k})|^2 \rangle_{|\bm{k}|=k}$ and the scaled spectrum $D(k) = k^2 P(k)$.
For HEALPix maps, we compute the angular power spectrum $C_\ell$ via \texttt{healpy.anafast} and form $D_\ell = \ell(\ell+1)C_\ell / 2\pi$.

\paragraph{Gaussianity metrics.}
For each temperature field, we compute the skewness $\gamma_1 = \langle T^3 \rangle / \sigma^3$ and excess kurtosis $\gamma_2 = \langle T^4 \rangle / \sigma^4 - 3$, with theoretical standard errors $\sigma_{\gamma_1} = \sqrt{6/N_{\mathrm{pix}}}$ and $\sigma_{\gamma_2} = \sqrt{24/N_{\mathrm{pix}}}$.
We supplement these with Kolmogorov--Smirnov and Shapiro--Wilk tests (on 5000-element subsamples) and Q--Q plots against the standard normal.

\paragraph{Spectral correlation.}
To compare DSC and $\Lambda$CDM spectral shapes, we compute the Pearson correlation coefficient between $\ln D(k)_{\mathrm{DSC}}$ and $\ln D(k)_{\Lambda\mathrm{CDM}}$ over the range $k \in [2, 60]$, after normalizing both to equal integrated power.

\paragraph{Toy analytic reference spectrum.}
For 2D power spectrum shape comparisons, we use a simplified analytic formula
\begin{equation}
  D(k)_{\mathrm{ref}} = e^{-(k/35)^2}\!\left[1 + 1.2\sin^2\!\left(\frac{\pi k}{12}\right)\right]\frac{50}{k + 10}\,,
  \label{eq:mock}
\end{equation}
which captures the qualitative features of acoustic oscillation peaks and Silk damping in a simplified form.
We stress that this is a toy reference used only for qualitative shape comparison; quantitative benchmarking against the observed CMB is performed in \cref{sec:chi2} using the real Planck SMICA angular power spectrum~\citep{Planck2018V}.

\subsection{Simulation-based inference}
\label{sec:sbi}

To test whether the DSC lattice parameters can be inferred from observed power spectra, we perform a twin experiment using Bayesian optimization~\citep{Akiba2019}.
A ``ground truth'' universe is simulated with hidden parameters $(c^2 = 0.650, \eta = 0.015, \kopt = 2.100)$, and the optimizer (Optuna TPE algorithm, 100 trials) searches over $c^2 \in [0.3, 0.9]$, $\eta \in [0.005, 0.05]$, $\kopt \in [1.0, 4.0]$ to minimize the mean squared error between normalized $D(k)$ spectra.
The initial field is fixed (seed = 42) to isolate parameter effects.
